%% ma: L12-B19-0032
%% lop: 12
%% chuong: 6
%% bai: 19
%% dang: 12.19.12
%% loai: TN4
%% muc: NB
%% dap_an: "C"
%% nguon: "(NBV)-10 ĐỀ KIỂM TRA HK2-TOÁN 12-DUNG CHUNG 3 SÁCH.pdf (H03, Đề số 3), Phần 1 Câu 11"
%% kien_thuc: "Bài 19 (công thức xác suất toàn phần)"
%% trang_thai: cho-duyet
%% ly_do: "Dạng toán mới đề xuất 12.19.12, chờ giáo viên duyệt"
%% ghi_chu: "Đáp án gốc C đúng."
\begin{ex}
Cho các biến cố $A$ và $B$ sao cho $0<P(A)<1$, $0<P(B)<1$. Khẳng định nào sau đây là đúng?
\choice{$P(B)=P(B)\cdot P(A\mid B)+P\left(\overline B\right)\cdot P\left(A\mid\overline B\right)$}{$P(A)=P(B)\cdot P(A\mid B)+P\left(\overline B\right)\cdot P\left(\overline A\mid B\right)$}{\True $P(A)=P(B)\cdot P(A\mid B)+P\left(\overline B\right)\cdot P\left(A\mid\overline B\right)$}{$P(B)=P(A)\cdot P(B\mid A)+P(B)\cdot P(A\mid B)$}
\loigiai{Theo công thức xác suất toàn phần với hệ $\{B,\overline B\}$: $P(A)=P(B)\cdot P(A\mid B)+P\left(\overline B\right)\cdot P\left(A\mid\overline B\right)$. Chọn C.}
\end{ex}
